\documentclass[11pt,letterpaper]{article}
\usepackage[margin=0.82in,headheight=14pt,headsep=16pt,footskip=25pt]{geometry}
\usepackage[T1]{fontenc}
\usepackage{lmodern,amsmath,amssymb,microtype,enumitem,xcolor,titlesec,fancyhdr,hyperref}
\definecolor{accent}{HTML}{21566D}
\hypersetup{colorlinks=true,urlcolor=accent,linkcolor=accent,pdftitle={SciCode 2: Quantum Error-Correction Decoding on D-Wave},pdfauthor={SciCode2}}
\pagestyle{fancy}\fancyhf{}
\fancyhead[L]{\footnotesize\sffamily SciCode 2 | Decoding formulation and embedding}
\fancyhead[R]{\footnotesize\sffamily Problem proposal}
\fancyfoot[C]{\small\thepage}
\renewcommand{\headrulewidth}{0.3pt}
\setlength{\parindent}{0pt}\setlength{\parskip}{4pt}
\setlist{leftmargin=*,itemsep=3pt,topsep=3pt,parsep=0pt}
\titleformat{\section}{\large\bfseries\color{accent}}{\thesection}{0.55em}{}
\titlespacing*{\section}{0pt}{10pt}{5pt}
\setlength{\emergencystretch}{2em}
\begin{document}
\thispagestyle{plain}
{\color{accent}\small\sffamily SCI CODE 2\quad|\quad RESEARCH BENCHMARK}\par\vspace{8pt}
{\LARGE\bfseries Quantum Error-Correction\\[2pt] Decoding on D-Wave}\par\vspace{5pt}
{\large Formulation and embedding on Advantage2}\par\vspace{5pt}
{\small Problem proposal\quad\textbullet\quad 7 October 2026}\par\vspace{7pt}
{\color{accent}\hrule}\vspace{10pt}

\textbf{Scientific objective.} Given quantum-code check matrices and measured syndromes, construct a faithful quadratic annealing model and a resource-efficient embedding on D-Wave Advantage2 connectivity. Determine how the decoding representation and its embedding jointly affect physical resources and programmable energy scales. Work offline; QPU access is not required.\par

\section{Background}

Quantum error-correction decoding has been formulated for Ising hardware [4] and classical annealing systems [5,6]. The present question concerns compiling matrix-defined decoding problems onto Advantage2 connectivity before QPU access.\par

\section{Why this is unresolved}

A parity-check matrix does not directly specify the quadratic graph programmed on a quantum annealer. Auxiliary variables, equivalent constraints and the choice of decoding representation can change connectivity and coefficient scales. Embedding adds physical variables and couplings. A compact logical formulation may therefore have a costly embedding or a poor usable energy scale.\par

\section{Research question}

Construct and validate resource-efficient quadratic encodings and minor embeddings of check-matrix-defined quantum decoding problems on D-Wave Advantage2 connectivity, and determine which formulation and embedding choices remain favorable under hardware constraints.\par

Seek a justified resource frontier within a declared classical search budget. Global embedding optimality is not required; reserve “optimal” for a specified objective with a certificate. The scientific choices of representation, reduction, placement, penalty strengths and recovery remain open.\par

\section{Mathematical problem and inputs}

All code relations use binary arithmetic over $\mathbb F_2$. The required core is one error sector of a CSS stabilizer code, specified by $H_X\in\mathbb F_2^{m_X\times n}$ and $H_Z\in\mathbb F_2^{m_Z\times n}$ with $H_XH_Z^T=0$. For X errors set $H=H_Z$ and $G=H_X$; exchange them for Z errors. The rows of $G$ generate stabilizer-equivalent corrections. A lone $H$ specifies syndrome constraints but does not determine quantum logical equivalence.\par

An instance contains a syndrome $s\in\operatorname{im}(H)$ and positive weights $w_i$. Its reference optimization is\par

\[C^*(H,s,w)=\min_{e\in\{0,1\}^n}\sum_{i=1}^n w_i e_i\quad\text{subject to}\quad He=s\pmod 2.\]

For independent errors with $0<p_i<1/2$, use $w_i=\log[(1-p_i)/p_i]$; uniform weights give minimum-weight decoding. Choose and report noise probabilities, syndrome suites and seeds. A suite can be generated by independently drawing hidden errors $e_{\rm true}$ and forming $s=He_{\rm true}$; the decoder receives only the matrices, syndrome and weights. Declare the tie rule. Minimum-cost error decoding is the required optimization; degenerate maximum-likelihood decoding, which sums probability within logical classes, is an optional extension.\par

Derive a linear logical-label map $L$ complete on $\ker H$, meaning $\ker(L|_{\ker H})=\operatorname{row}(G)$. A correction succeeds on a hidden error precisely when\par

\[H\hat e=s,\qquad \hat e\oplus e_{\rm true}\in\operatorname{row}(G).\]

For a syndrome-valid correction the second condition is equivalent to $L(\hat e\oplus e_{\rm true})=0$. Count an invalid correction or missing output as failure. Exact error-string agreement is not the success criterion. Hidden errors and their logical classes are available only to evaluation, never to the decoder or its preprocessing.\par

\section{Reproducible code families}

Use a periodic square-lattice toric family [8] with integer $d\geq3$ and $n=2d^2$. For $x,y\in\{0,\ldots,d-1\}$, index horizontal and vertical edges by $h(x,y)=2(x+dy)$ and $v(x,y)=h(x,y)+1$. Reduce every coordinate modulo $d$ before indexing. Index check rows by $x+dy$. The support of a vertex row of $H_X$ is\par

\[\{h(x,y),h(x-1,y),v(x,y),v(x,y-1)\},\]

and the support of a face row of $H_Z$ is\par

\[\{h(x,y),v(x+1,y),h(x,y+1),v(x,y)\}.\]

Set these entries to one and all others to zero. Keep all wraparound checks; redundant rows are allowed. Start with $d=3$ and choose larger sizes to probe resource growth. Derive the logical labels from the matrices rather than assuming a planar boundary. This construction defines the full input without a separate dataset.\par

Include a non-toric CSS example through the same interface: the seven-qubit Steane code [10], with\par

\[H_X=H_Z=\begin{pmatrix}1&0&1&0&1&0&1\\0&1&1&0&0&1&1\\0&0&0&1&1&1&1\end{pmatrix}.\]

Use this as a small transfer and correctness case; it need not support a scaling claim. Additional non-toric codes are welcome if their matrices or public deterministic construction are supplied. For every example, report matrix ranks and verify commutation. Geometry may be used when available, but the code must remain defined by its matrices.\par

\section{Target graph and hardware contract}

Advantage2 uses Zephyr connectivity; earlier Advantage systems use Pegasus [1]. The offline core uses ideal Zephyr graphs generated by the public Ocean function \texttt{dwave.graphs.zephyr\_graph(m, t=4)} [9], with a positive integer grid parameter $m$ chosen and reported for each comparison. This graph has $16m(2m+1)$ vertices and maximum degree at most 20. Save the generator/package version, arguments, node and edge lists, and hashes. Compare alternatives on identical target graphs; report target-size growth separately from embedding improvements.\par

This public generator is sufficient to start without credentials or a device export. A graph with deleted nodes or couplers is an optional robustness case; an exported QPU working graph is an optional device-specific case. Label ideal, synthetic-defect and actual device graphs separately. Toric periodicity belongs to the code, not to the hardware: do not add wraparound hardware edges absent from the chosen Zephyr graph.\par

Device-specific programmability claims require a solver identifier, graph revision and documented bias, coupling and coupling-group limits [2]. Without those limits, give a parameterized coefficient-scale analysis and conditional hardware feasibility. An ideal graph alone establishes topology compatibility, not availability on a particular QPU.\par

\section{Encoding and embedding requirements}

Construct a quadratic model\par

\[Q(y)=c+\sum_u a_u y_u+\sum_{u<v}b_{uv}y_uy_v,\qquad y_u\in\{0,1\},\]

with a documented projection $\pi(y)=e$. Prove that every projected ground state is a feasible minimum-cost correction and that the reference optimum is attained. A justified tie rule need not represent every tied minimizer. Do not omit a logical sector without showing that the original optimum is preserved. Document auxiliary variables, penalty conditions, offsets and positive scale factors; a higher-order parity expression alone is not a QUBO.\par

Embed the actual interaction graph, including isolated variables. Each logical variable must map to a nonempty connected branch set in the target, branch sets must be pairwise disjoint, and every nonzero logical interaction must have a physical coupler between its branch sets. Report physical biases, couplings, chain constraints and the conversion to a stated Ising convention. For consistent branch sets, physical energy must reproduce logical energy up to a known positive scale and constant. Justify any stronger assertion about physical ground states. Include penalties and chain couplings when assessing coefficient limits and usable scale.\par

Charge all preprocessing, reduction, embedding search and sample recovery to the relevant cost. In particular, if syndrome-dependent classical preprocessing already solves decoding, identify the result as classical decoding followed by answer encoding. Encoding a previously computed answer does not by itself establish an improved annealing compiler. Report both initial and amortized costs when reusing an embedding across syndromes.\par

\section{Scope and approach}

Use perfect syndrome measurements and independent binary noise in the required CSS sector. Select tractable sizes and declare time, memory, search and sampling budgets before comparing alternatives. Public literature, public software and internet access are allowed. No private source, proprietary dataset, repository history or QPU account is needed. Correlated noise, general stabilizer codes and noisy syndrome histories are optional extensions.\par

\section{Evidence requirements}

\begin{itemize}\item \textbf{Correctness:} Independently verify code relations, quadratic reduction, projection, embedding connectivity and physical coefficient construction. Exact small cases must check the ground-state projection, ties and auxiliary states; label uncertified larger cases heuristic.
\item \textbf{Controlled resource comparison:} Compare at least two defensible formulation/embedding designs, including a reproducible automatic Ocean/minorminer reference [3,7]. Control the logical model when comparing embeddings and the target graph when comparing formulations. Report logical variables/interactions, physical qubits/couplers, branch-set size distributions, runtime, peak memory, coefficient scale and failures. Explain how the evidence supports the claimed tradeoff.
\item \textbf{Transfer and reuse:} Include the toric family and Steane example, independent syndrome sets, complete instance specifications and random seeds. State which quantities change with syndrome and when an embedding can be reused.
\item \textbf{Recovery interface:} Specify how broken-chain or excited samples would be converted into corrections. If claiming a recovery benefit, support it with controlled offline samples, runtime and before/after syndrome and logical outcomes. Artificial perturbations or classical annealer samples are diagnostics, not QPU error distributions.
\item \textbf{Decoding claims:} If assessing decoding quality, use identical information, noise, syndromes and declared budgets for an appropriate classical decoder and a classical optimizer on the same quadratic model. Separate logical failure with uncertainty, syndrome feasibility and optimization success; include recovery and restart costs.\end{itemize}

A valid conclusion may be a favorable construction, a tradeoff without a universal winner, or a reproducible capacity limitation. A failed heuristic search does not prove non-embeddability, and a local minimum need not be a correct decoding solution. Offline results cannot establish quantum speedup, QPU sampling performance or the best physical annealing parameters. No hardware run, threshold estimate or real-time latency claim is required.\par

\section{Deliverables and acceptance}

\begin{enumerate}\item A documented matrix-to-QUBO construction with mathematical justification and reusable source code, plus reproducible code, syndrome and target-graph generators.
\item Portable models, projections, embeddings, hardware manifests and independent validation results, including a boundary-correct toric example suitable for external technical review.
\item Machine-readable correctness and resource comparisons, figures of the resource tradeoffs, and an English report explaining decisions, evidence and limitations. Provide dependency versions, commands and seeds so another researcher can reproduce the work in a fresh directory.\end{enumerate}

Success means a faithful, independently validated construction and a controlled, decision-relevant comparison on the declared instance suites, target graphs and budgets. A negative result needs the same reproducible evidence; a single automatic embedding or a timeout alone is insufficient. Methods and file layout are open. No prescribed algorithm or unreported numerical performance threshold determines acceptance.\par

\section{Expected outcomes}

A smaller quadratic model may have a larger physical footprint or worse coefficient scale. Different representations and placements may favor different sizes or code families. Embedding reuse may reduce amortized cost if the representation retains a compatible interaction graph. These are hypotheses to investigate, not required positive results.\par

\clearpage

\section{Public references}

{\small [1] D-Wave, \href{https://docs.dwavequantum.com/en/latest/quantum_research/topologies.html}{Topologies}.\par}\vspace{5pt}

{\small [2] D-Wave, \href{https://docs.dwavequantum.com/en/latest/quantum_research/solver_properties_all.html}{QPU Solver Properties}.\par}\vspace{5pt}

{\small [3] D-Wave, \href{https://docs.dwavequantum.com/en/latest/ocean/api_ref_system/composites.html}{Ocean embedding composites}.\par}\vspace{5pt}

{\small [4] J. Roffe et al., \href{https://arxiv.org/abs/1903.10254}{Decoding quantum error correction with Ising model hardware}, arXiv:1903.10254 (2019).\par}\vspace{5pt}

{\small [5] J. Fujisaki et al., \href{https://arxiv.org/abs/2203.15304}{A Practical and Scalable Decoder for Topological Quantum Error Correction with Digital Annealer}, Phys. Rev. Research \textbf{4}, 043086 (2022).\par}\vspace{5pt}

{\small [6] Y. Takeuchi et al., \href{https://arxiv.org/abs/2311.07973}{Comparative study of decoding the surface code using simulated annealing under depolarizing noise}, arXiv:2311.07973v2 (2023).\par}\vspace{5pt}

{\small [7] J. Cai, W. G. Macready and A. Roy, \href{https://arxiv.org/abs/1406.2741}{A practical heuristic for finding graph minors}, arXiv:1406.2741 (2014).\par}\vspace{5pt}

{\small [8] E. Dennis et al., \href{https://arxiv.org/abs/quant-ph/0110143}{Topological quantum memory}, J. Math. Phys. \textbf{43}, 4452 (2002).\par}\vspace{5pt}

{\small [9] D-Wave, \href{https://docs.dwavequantum.com/en/latest/ocean/api_ref_graphs/generated/dwave.graphs.zephyr_graph.html}{Zephyr graph generator}. Public software documentation accessed 7 October 2026.\par}\vspace{5pt}

{\small [10] A. Steane, \href{https://arxiv.org/abs/quant-ph/9601029}{Multiple Particle Interference and Quantum Error Correction}, Proc. R. Soc. Lond. A \textbf{452}, 2551 (1996).
\par}\vspace{5pt}
\end{document}
